% 2021 全国甲卷文科 2：农户家庭年收入的频率分布直方图
% 绘制依据：源 PDF 第 1 页原图。纵轴「频率/组距」，横轴「收入/万元」。
% 横轴刻度 0 与 2.5,3.5,…,14.5；纵轴标注 0.02、0.04、0.10、0.14、0.20。
% 十二个组距均为 1，自 2.5 起各组高度依次为
%   0.02, 0.04, 0.10, 0.14, 0.20, 0.20, 0.10, 0.10, 0.04, 0.02, 0.02, 0.02。
% 交叉验算（四项独立约束全部吻合，故读数可靠）：
%   Σ高度×1 = 1.00；低于 4.5 者 0.02+0.04 = 6%（选项 A）；
%   不低于 10.5 者 0.04+0.02×3 = 10%（选项 B）；4.5~8.5 者 0.64 > 0.5（选项 D）；
%   均值 Σ组中值×频率 = 7.68 > 6.5（选项 C 故为「不正确」者）。
% 版式：横轴 0 与首刻度 2.5 之间画断轴锯齿；每个纵轴标注值处引水平虚线至对应柱顶。
\begin{tikzpicture}[font=\small]
  % x: 1 单位收入 = 0.62cm，自 2.5 起；y: 0.02 = 0.42cm
  \def\xs{0.62}\def\ys{21.0}
  \pgfmathsetmacro{\xz}{0.55}% 断轴后横轴起点（对应收入 2.5）
  % 断轴的横轴
  \draw (0,0) -- (0.10,0);
  \draw (0.10,0) -- (0.16,0.10) -- (0.26,-0.10) -- (0.32,0);
  \draw[->] (0.32,0) -- ({\xz+12*\xs+0.75},0);
  \node[anchor=north] at (0,-0.06) {$0$};
  \node[anchor=west] at ({\xz+12*\xs+0.30},0.26) {收入$/$万元};
  % 纵轴
  \draw[->] (0,0) -- (0,{0.20*\ys+0.85});
  \node[anchor=south east, align=right] at (0.06,{0.20*\ys+0.30}) {频率\\[-2pt]\rule{1.1cm}{0.4pt}\\[-2pt]组距};
  % 纵轴标注与引至柱顶的水平虚线
  \foreach \v/\xend in {0.02/12, 0.04/9, 0.10/7, 0.14/4, 0.20/3} {
    \node[anchor=east, inner sep=2pt] at (-0.04,{\v*\ys}) {$\v$};
    \draw[dashed, gray] (0,{\v*\ys}) -- ({\xz+\xend*\xs},{\v*\ys});
  }
  % 十二个矩形
  \foreach \i/\h in {0/0.02, 1/0.04, 2/0.10, 3/0.14, 4/0.20, 5/0.20,
                     6/0.10, 7/0.10, 8/0.04, 9/0.02, 10/0.02, 11/0.02} {
    \draw ({\xz+\i*\xs},0) rectangle ({\xz+(\i+1)*\xs},{\h*\ys});
  }
  % 横轴刻度标注
  \foreach \i/\t in {0/2.5, 1/3.5, 2/4.5, 3/5.5, 4/6.5, 5/7.5, 6/8.5,
                     7/9.5, 8/10.5, 9/11.5, 10/12.5, 11/13.5, 12/14.5} {
    \node[anchor=north, inner sep=2pt, font=\tiny] at ({\xz+\i*\xs},-0.06) {$\t$};
  }
\end{tikzpicture}
